\section{Rolling hash}
\whenuse{Comparar substrings, detectar repeticiones o hacer binary search sobre
longitudes aceptando una probabilidad peque\~na de colisi\'on.}

Hash doble para comparar substrings en $O(1)$ tras preprocesar. Los rangos son
semiabiertos $[l,r)$. Una igualdad de hashes sigue siendo probabil\'istica.

\begin{lstlisting}[style=cpp]
struct RollingHash {
    static const int MOD1 = 1000000007;
    static const int MOD2 = 1000000009;
    static const int BASE = 911382323;
    vector<int> power1, power2, hash1, hash2;

    explicit RollingHash(const string &s) {
        int n = s.size();
        power1.assign(n + 1, 1); power2.assign(n + 1, 1);
        hash1.assign(n + 1, 0);  hash2.assign(n + 1, 0);
        for (int i = 0; i < n; ++i) {
            int value = (unsigned char)s[i] + 1;
            power1[i + 1] = (ll)power1[i] * BASE % MOD1;
            power2[i + 1] = (ll)power2[i] * BASE % MOD2;
            hash1[i + 1] = ((ll)hash1[i] * BASE + value) % MOD1;
            hash2[i + 1] = ((ll)hash2[i] * BASE + value) % MOD2;
        }
    }

    pair<int,int> get(int l, int r) const {
        ll first = hash1[r] - (ll)hash1[l] * power1[r - l] % MOD1;
        ll second = hash2[r] - (ll)hash2[l] * power2[r - l] % MOD2;
        if (first < 0) first += MOD1;
        if (second < 0) second += MOD2;
        return {(int)first, (int)second};
    }
};
\end{lstlisting}

\complexity{construcci\'on $O(n)$; consulta $O(1)$; memoria $O(n)$.}

Comparar dos substrings de la misma longitud:

\begin{lstlisting}[style=cpp]
string s="banana";
RollingHash hash(s);
bool equal=hash.get(1,4)==hash.get(3,6);
cout<<equal<<'\n'; // 1: "ana" == "ana"
\end{lstlisting}

Buscar apariciones de un patrón:

\begin{lstlisting}[style=cpp]
string text="bananana",pattern="ana";
RollingHash text_hash(text),pattern_hash(pattern);
auto target=pattern_hash.get(0,pattern.size());
for(int i=0;i+(int)pattern.size()<=(int)text.size();++i)
    if(text_hash.get(i,i+pattern.size())==target)
        cout<<i<<' '; // 1 3 5
\end{lstlisting}

\texttt{get(l,r)} devuelve el hash del substring semiabierto $[l,r)$. Sólo
deben compararse substrings de igual longitud. Aunque se use hash doble, existe
una probabilidad muy pequeña de colisión.